\documentclass{csmathnotes}
\usepackage{csquotes}
\usepackage{biblatex}
\usepackage[linesnumbered,lined,ruled]{algorithm2e}
\usepackage[dvipsnames]{xcolor}
\usepackage{adjustbox}
\usepackage{booktabs}
\usepackage{array}
\setcounter{biburlnumpenalty}{9000}
\setcounter{biburllcpenalty}{9000}
\setcounter{biburlucpenalty}{9000}

\udc{519.688}

\title{Сравнение алгоритмов программного соперника в игровом приложении <<Коридор>>}

\author{Карпунин Александр Олегович}
\affiliation{Ярославский государственный университет им. П.Г. Демидова}
\email{aleksandrkarpunin0889@gmail.com}

\advisor{д. ф.-м. н., доцент}{Николаев Андрей Валерьевич}
\affiliation{Ярославский государственный университет им. П.Г. Демидова}
\advisoremail{andrei.v.nikolaev@gmail.com}

\addbibresource{bibliography.bib}

\begin{document}

\SetAlgorithmName{Алгоритм}{}
\SetAlFnt{\scriptsize}
\SetAlCapFnt{\scriptsize}
\SetAlCapNameFnt{\scriptsize}

\maketitle

\begin{abstract}
В статье рассматривается разработка модуля стратегического поведения программного соперника для настольной игры <<Коридор>>. Описана игровая модель, функция оценки позиции и три алгоритма выбора хода: MiniMax, AlphaBeta и MonteCarlo. Для оценки качества реализованных подходов проведено автоматическое сравнение алгоритмов на поле $11 \times 11$. В результате эксперимента AlphaBeta демонстрирует наибольшую долю побед, а наибольшее время хода необходимо MonteCarlo.

\keywords{игра <<Коридор>>, программный соперник, алгоритмы поиска, MonteCarlo, функция оценки}
\end{abstract}

\section{Введение}

Логические настольные игры удобны для исследования алгоритмов искусственного интеллекта, поскольку состояние партии можно строго описать, а множество допустимых действий формируется по заранее заданным правилам. Одной из таких игр является <<Коридор>> "--- игра с полной информацией для двух игроков. Каждый игрок управляет фишкой и стремится первым дойти до противоположной стороны поля. В процессе игры игроки могут устанавливать перегородки, усложняющие путь соперника~\cite{QuoridorRules}.

В <<Коридоре>> каждый игрок выбирает между продвижением к цели и изменением структуры игрового поля. Поэтому программный соперник должен учитывать не только собственное расстояние до финишной линии, но и возможность замедлить соперника с помощью стен. Полный перебор дерева игры быстро становится слишком трудозатратным, поэтому алгоритмы выбора хода требуют ограничений глубины, эвристик и оценки промежуточных позиций.

Цель работы "--- разработать модуль программного соперника для игрового приложения <<Коридор>> и сравнить несколько алгоритмов выбора хода: MiniMax, AlphaBeta и MonteCarlo.

\section{Правила и модель игры}

В настольной игре используется квадратное поле $9 \times 9$, однако в разработанном приложении также поддерживаются поля $7 \times 7$ и $11 \times 11$. В начале партии фишки игроков располагаются на противоположных сторонах поля: первый игрок "--- в центре нижней строки, второй "--- в центре верхней. Цель каждого игрока "--- первым достичь любой клетки своей целевой линии.

\begin{figure}
    \centering
    \includegraphics[width=0.65\textwidth]{images/game-board.png}
    \caption{Игровое поле приложения <<Коридор>>}
    \label{fig:game_board}
\end{figure}

За один ход игрок может выполнить одно из двух действий~\cite{QuoridorRules}:
\begin{itemize}
    \item Передвинуть фишку. Фишка передвигается на смежную клетку, если переход не заблокирован стеной. Если соседняя клетка занята фишкой соперника и за ней нет ни перегородки, ни границы поля, фишка может перепрыгнуть через соперника. Если прямой прыжок невозможен, допускается прыжок на клетку слева или справа от той, куда изначально предполагался прыжок.
    \item Установить стену. Стена занимает пространство между клетками, имеет длину две клетки и может быть расположена горизонтально или вертикально. Установка стены допускается, если у игрока ещё остались перегородки, новая стена не пересекается с уже установленными и после её установки у каждого игрока остаётся хотя бы один маршрут к своей целевой линии.
\end{itemize}

Для проверки существования маршрута используется поиск в ширину от текущей позиции фишки до ближайшей клетки целевой строки. Если после временного применения стены путь хотя бы одного игрока исчезает, такой ход исключается из множества допустимых.

В модуле стратегического поведения множество допустимых ходов формируется из двух частей. Первая часть "--- перемещения фишки, получаемые проверкой направлений хода с текущей клетки и обработкой возможных прыжков через соперника. Вторая часть "--- варианты установки стен около кратчайшего пути соперника. Такой выбор используется потому, что стена полезна прежде всего тогда, когда меняет текущий лучший маршрут соперника; перегородка вдали от этого маршрута часто не увеличивает длину пути и только расширяет перебор.

Состояние игры можно представить набором данных:
\begin{itemize}
    \item расположение фишек игроков;
    \item расположение стен;
    \item количество оставшихся стен у каждого игрока;
    \item номер игрока, который должен сделать следующий ход.
\end{itemize}

Такое представление позволяет использовать классические алгоритмы поиска в дереве игры~\cite{Russell:2021}. Узел дерева соответствует состоянию партии, а ребро "--- допустимому ходу. Лист дерева оценивается численно с помощью функции оценки позиции.

\section{Функция оценки позиции}

Поскольку полный перебор партии невозможен, алгоритмам требуется приближённая оценка промежуточного состояния. В разработанном приложении оценка позиции строится как линейная комбинация нескольких признаков:
\[
    F(s) = \sum_{k=1}^{7} \lambda_k f_k(s),
\]
где $s$ "--- состояние игры, $f_k(s)$ "--- числовые признаки игровой позиции, а $\lambda_k$ "--- соответствующие им веса. Чем больше итоговое значение $F(s)$, тем выгоднее позиция для программного соперника.

Основной признак связан с кратчайшими путями игроков до целевых строк. Если путь алгоритма до цели короче пути соперника, позиция считается более выгодной. Дополнительно учитываются мобильность фишек, оставшиеся стены, продвижение к целевой строке и специальные эндшпильные признаки, которые усиливают оценку позиций, близких к победе или поражению. Используемые признаки приведены в таблице~\ref{tab:features}.

\begin{table}[h]
\centering
\caption{Признаки функции оценки позиции}
\label{tab:features}
\resizebox{\textwidth}{!}{%
\begin{tabular}{clp{8cm}}
\toprule
$k$ & Признак $f_k(s)$ & Смысл признака \\
\midrule
1 & $d_p - d_a$ & Разность кратчайших путей соперника и алгоритма до целевых строк \\
2 & $\mathrm{mob}_a$ & Количество допустимых перемещений фишкой алгоритма \\
3 & $-\mathrm{mob}_p$ & Штраф за количество допустимых перемещений фишкой соперника \\
4 & $w_a - w_p$ & Преимущество алгоритма по числу оставшихся стен \\
5 & $\mathrm{prog}_a - \mathrm{prog}_p$ & Разность продвижения игроков к своим целевым строкам \\
6 & $e(d_a)$ & Бонус за близость алгоритма к победе \\
7 & $-e(d_p)$ & Штраф за близость соперника к победе \\
\bottomrule
\end{tabular}%
}
\end{table}

Здесь $d_a$ и $d_p$ "--- длины кратчайших путей алгоритма и соперника, $w_a$ и $w_p$ "--- количество оставшихся стен. Функция $e(d)$ задаёт эндшпильную поправку:
\[
e(d) =
\begin{cases}
10, & d \le 1, \\
4,  & d = 2, \\
1,  & d = 3, \\
0,  & d > 3.
\end{cases}
\]
То есть позиция, в которой алгоритму остаётся один ход до победы, получает большой положительный вклад, а позиция, где близок к победе соперник, получает симметричный штраф.

Начальные веса были заданы вручную на основе смысла признаков:
\[
\lambda_1 = 65,\quad
\lambda_2 = 4,\quad
\lambda_3 = 3,\quad
\lambda_4 = 18,\quad
\lambda_5 = 6,\quad
\lambda_6 = 165,\quad
\lambda_7 = 175.
\]
Наибольшие веса имеют эндшпильные признаки и разность кратчайших путей, поскольку для игры <<Коридор>> особенно важно не только продвигать собственную фишку, но и вовремя замечать угрозу быстрого выхода соперника к целевой строке.

При оценке стен дополнительно используется локальная эвристика стратегической ценности. После временного применения стены на копии поля вычисляется, насколько изменились кратчайшие пути игроков:
\[
\mathrm{wallImpact} =
(d_p^{\mathrm{after}} - d_p^{\mathrm{before}}) \cdot 100
- \max(0, d_a^{\mathrm{after}} - d_a^{\mathrm{before}}) \cdot 45.
\]
Таким образом, стена получает высокий приоритет, если она увеличивает путь соперника, и штрафуется, если одновременно ухудшает путь самого алгоритма. Это используется при сортировке ходов и при поиске защитного хода в эндшпиле.

В приложении также реализован механизм коррекции весов по результатам self-play партий. Для каждого хода алгоритма сохраняется обучающий образец: значения признаков до хода и после него. После завершения партии вычисляется изменение каждого признака $\Delta f_k = f_k^{\mathrm{after}} - f_k^{\mathrm{before}}$, после чего формируется суммарный сигнал:
\[
g_k = \sum_{\text{ходы алгоритма}} \mathrm{reward} \cdot \Delta f_k,
\quad
\mathrm{reward} =
\begin{cases}
+1, & \text{алгоритм победил}, \\
-1, & \text{алгоритм проиграл}.
\end{cases}
\]
Если $g_k > 0$, увеличение соответствующего признака чаще встречалось в победных партиях, поэтому вес признака увеличивается на один шаг. Если $g_k < 0$, вес уменьшается. Новое значение дополнительно ограничивается минимальной и максимальной границей, чтобы один признак не начал подавлять остальные:
\[
\lambda_k^{\mathrm{new}} =
\begin{cases}
\lambda_k^{\min}, & \lambda_k' < \lambda_k^{\min}, \\
\lambda_k^{\max}, & \lambda_k' > \lambda_k^{\max}, \\
\lambda_k', & \lambda_k^{\min} \le \lambda_k' \le \lambda_k^{\max}.
\end{cases}
\]
После обучения веса сохраняются в файл конфигурации и используются при следующих запусках приложения. На момент проведения экспериментов сохранённые веса имели вид:
\[
\begin{aligned}
\lambda_1 &= 120,\quad
\lambda_2 = 7,\quad
\lambda_3 = 12,\quad
\lambda_4 = 8,\\
\lambda_5 &= 25,\quad
\lambda_6 = 186,\quad
\lambda_7 = 170.
\end{aligned}
\]
По сравнению с начальными значениями сильнее всего выросли веса разности путей, ограничения мобильности соперника и продвижения к целевой строке. Это соответствует наблюдаемой стратегии игры: алгоритм чаще выбирает позиции, где он не просто приближается к победе сам, но и вынуждает соперника тратить дополнительные ходы на обход стен.

\section{Алгоритм MiniMax}

MiniMax действует по следующему принципу: алгоритм выбирает ход, максимизирующий оценку позиции, а для игрока "--- минимизирующий~\cite{Russell:2021}. Поиск строится до заданной глубины, после чего листья дерева оцениваются с применением функции оценки $F(s)$.

Недостаток MiniMax заключается в том, что число просматриваемых состояний растёт экспоненциально с глубиной поиска. Поэтому в приложении используется ограничение глубины и предварительная сортировка ходов: сначала рассматриваются перемещения, уменьшающие путь к цели, и стены, увеличивающие путь соперника.

В реализации алгоритм работает с итеративным углублением: сначала строится дерево небольшой глубины, затем глубина постепенно увеличивается, пока не достигнут заданный предел или не истёк лимит времени. Такой подход позволяет вернуть лучший найденный ход даже в том случае, если полный просмотр очередного уровня не был завершён. Для повторяющихся состояний используется кэш оценок, ключом которого является состояние доски, глубина и тип узла. Это снижает число повторных вычислений в ситуациях, когда одна и та же позиция достигается разными последовательностями ходов.

Перед рекурсивным просмотром множество допустимых ходов сортируется. В начало списка попадают немедленные победные ходы, затем перемещения, улучшающие кратчайший путь алгоритма, и стены с высоким значением \texttt{wallImpact}. После этого число рассматриваемых кандидатов ограничивается в соответствии с уровнем сложности. Поэтому MiniMax в приложении является не полным перебором всех ходов до заданной глубины, а эвристически направленным поиском по наиболее перспективной части дерева. Общая схема выбора оценки приведена в алгоритме~\ref{alg:minimax}.

\begin{algorithm}[!t]
\caption{MiniMax с итеративным углублением}\label{alg:minimax}
\SetKwFunction{FMM}{minimax}
\SetKwFunction{FApply}{apply}
\SetKwFunction{FEval}{evaluate}
\SetKwProg{Fn}{Функция}{:}{}
\Fn{\FMM{$s$, $d$, isMax}}{
    \If{в cache есть значение для $(s, d, isMax)$}{
        \KwRet{значение из cache}\;
    }
    \If{$\mathrm{terminal}(s)$}{
        \KwRet{терминальную оценку позиции}\;
    }
    \If{$d = 0$}{
        \KwRet{\FEval{$s$}}\;
    }
    $p \leftarrow$ игрок, которому принадлежит ход в узле $s$\;
    $M \leftarrow$ \texttt{getMoves}$(s, p)$\;
    \eIf{$isMax$}{
        $best \leftarrow -\infty$\;
        \ForEach{$m \in M$}{
            $best \leftarrow \max(best,\; \FMM{\FApply{$s$, $m$}, d - 1, false})$\;
        }
    }{
        $best \leftarrow +\infty$\;
        \ForEach{$m \in M$}{
            $best \leftarrow \min(best,\; \FMM{\FApply{$s$, $m$}, d - 1, true})$\;
        }
    }
    сохранить $best$ в cache для $(s, d, isMax)$\;
    \KwRet{$best$}\;
}
\KwRet{$m^*$}
\end{algorithm}

\section{Алгоритм AlphaBeta}

AlphaBeta "--- оптимизированная версия алгоритма MiniMax. Он вычисляет тот же результат при том же порядке листьев, но отбрасывает ветви, которые не могут повлиять на итоговый выбор~\cite{Russell:2021}. Для этого используются две границы: $\alpha$ "--- лучшая уже найденная оценка для максимизирующего игрока, и $\beta$ "--- лучшая оценка для минимизирующего игрока. Схема поиска с отсечениями показана в алгоритме~\ref{alg:alphabeta}.

\begin{algorithm}[!tbp]
\caption{AlphaBeta}\label{alg:alphabeta}
\SetKwFunction{FAB}{alphabeta}
\SetKwProg{Fn}{Функция}{:}{}
\Fn{\FAB{$s$, $d$, $\alpha$, $\beta$, isMax}}{
    \If{$\mathrm{terminal}(s)$}{
        \KwRet{терминальную оценку позиции}\;
    }
    \If{$d = 0$}{
        \eIf{isNoisyPosition$(s)$}{
            \KwRet{quiescence$(s, \alpha, \beta, isMax, 2)$}\;
        }{
            \KwRet{$F(s)$}\;
        }
    }
    $M \leftarrow$ \texttt{getMoves}$(s)$\;
    \eIf{$isMax$}{
        $\mathrm{best} \leftarrow -\infty$\;
        \ForEach{$m \in M$}{
            $v \leftarrow$ \FAB{$\mathrm{apply}(s, m)$, $d - 1$, $\alpha$, $\beta$, false}\;
            $\mathrm{best} \leftarrow \max(\mathrm{best}, v)$\;
            $\alpha \leftarrow \max(\alpha, \mathrm{best})$\;
            \If{$\beta \le \alpha$}{
                обновить killer-ходы и goodMoves\;
                \textbf{break}\tcp*{отсечение в MAX-узле}
            }
        }
    }{
        $\mathrm{best} \leftarrow +\infty$\;
        \ForEach{$m \in M$}{
            $v \leftarrow$ \FAB{$\mathrm{apply}(s, m)$, $d - 1$, $\alpha$, $\beta$, true}\;
            $\mathrm{best} \leftarrow \min(\mathrm{best}, v)$\;
            $\beta \leftarrow \min(\beta, \mathrm{best})$\;
            \If{$\beta \le \alpha$}{
                обновить killer-ходы и goodMoves\;
                \textbf{break}\tcp*{отсечение в MIN-узле}
            }
        }
    }
    \KwRet{$\mathrm{best}$}\;
}
\end{algorithm}

То, насколько AlphaBeta будет эффективно отсекать ветки, сильно зависит от порядка их рассмотрения. Чем раньше будут рассмотрены более перспективные ходы, тем чаще будут происходить отсечения. Поэтому в реализации используются эвристики сортировки и запоминание ходов, которые ранее приводили к отсечениям. Благодаря этому алгоритм успевает рассмотреть дерево на большую глубину, чем Minimax за то же выделенное время.

Для ускорения поиска AlphaBeta использует несколько дополнительных эвристик. Первая из них "--- \textit{killer-ходы}: если некоторый ход на данной глубине уже приводил к отсечению $\beta \le \alpha$, он считается потенциально сильным и в похожих узлах рассматривается раньше остальных. Вторая эвристика хранит накопленную статистику успешных отсечений для ходов в таблице \texttt{goodMoves}. Чем чаще ход вызывал отсечение, тем больший приоритет он получает при последующей сортировке.

Отдельная доработка связана с позициями около конца партии. Если на предельной глубине один из игроков находится не дальше двух ходов от цели, позиция считается острой. В таком случае алгоритм не сразу применяет статическую оценку $F(s)$, а дополнительно рассматривает немедленные победные ходы и наиболее ценные защитные стены. Это уменьшает вероятность ситуации, когда поиск остановился прямо перед решающим ходом и неверно оценил позицию как безопасную.

\section{Алгоритм MonteCarlo}

MonteCarlo-алгоритм выбирает ход на основе многократных симуляций продолжения партии. Он не стремится равномерно раскрыть дерево до фиксированной глубины, а оценивает ходы статистически. В работе~\cite{Respall:2018} для игры <<Коридор>> аналогично рассматривается алгоритм MCTS, что подтверждает применимость такого подхода к данной задаче. Общая схема статистического выбора хода приведена в алгоритме~\ref{alg:montecarlo}.

В разработанном приложении для каждого доступного хода выполняются симуляции, в которых дальнейшие действия выбираются с учётом функции оценки позиции. Чем чаще ход приводит к выигрышным или позиционно выгодным продолжениям, тем выше его итоговый рейтинг.

Поиск MonteCarlo организован как дерево посещённых состояний. Каждый узел хранит число посещений $n_v$ и число успешных симуляций $w_v$. При выборе следующего узла используется показатель UCT:
\[
\mathrm{UCT}(v) =
\frac{w_v}{n_v}
+ C \sqrt{\frac{\ln n_{\mathrm{parent}}}{n_v}}
+ h(v).
\]
Первое слагаемое отражает накопленную долю успешных симуляций, второе отвечает за исследование редко посещённых узлов, а $h(v)$ добавляет небольшую эвристическую поправку на основе функции оценки позиции. В реализации эта поправка вычисляется как $0{,}15 \cdot \tanh(F(s_v) / 500)$, поэтому она помогает направлять поиск, но не подавляет статистику симуляций.

Одна итерация MonteCarlo состоит из четырёх этапов: выбора узла по UCT, расширения дерева новым дочерним состоянием, симуляции партии и обратного распространения результата. Во время симуляции алгоритм не делает полностью случайные ходы: сначала проверяется немедленная победа, затем возможность поставить защитную стену, если соперник близок к цели, а в остальных случаях ход выбирается из ограниченной выборки по эвристической оценке. После завершения симуляции статистика посещений и результатов обновляется на всём пути от выбранного узла к корню.

\begin{algorithm}[!tbp]
\caption{MonteCarlo}\label{alg:montecarlo}
\KwIn{состояние $s$, временной лимит $T$, лимит итераций $N$}
\KwOut{выбранный ход}
$r \leftarrow$ корневой узел с состоянием $s$\;
\While{не истёк лимит $T$ \textbf{и} итераций $< N$}{
    $v \leftarrow \mathrm{select}(r)$ \tcp*{спуск по UCT}
    \If{$\neg\,\mathrm{terminal}(v.s)$}{
        $v \leftarrow \mathrm{expand}(v)$ \tcp*{новый дочерний узел}
    }
    $\mathrm{result} \leftarrow \mathrm{simulate}(v, \mathrm{rolloutDepth})$\;
    $\mathrm{backpropagate}(v, \mathrm{result})$\;
    \If{$r.n > 300$ \textbf{и} $\max_c(w_c/n_c) > 0.97$}{
        \textbf{break}\;
    }
}
\KwRet{$\arg\max_{c \in \mathrm{children}(r)} (n_c + w_c / n_c)$}
\end{algorithm}

Алгоритм MonteCarlo может работать без точного перебора дерева до заданной глубины, но у него есть существенный недостаток: для устойчивого выбора требуется провести много симуляций, что подразумевает значительные временные затраты.

\section{Методика сравнения}

Для сравнения алгоритмов в приложении реализован специальный режим проведения автоматических партий.

Эксперимент проводился на большом поле $11 \times 11$. Для каждой пары алгоритмов запускалось по 50 партий, роли первого и второго игрока чередовались. Для всех алгоритмов использовался максимальный уровень сложности. В режиме сравнения обучение весов функции оценки не выполнялось, чтобы параметры оставались фиксированными. В эксперименте также участвовал алгоритм случайного выбора хода.

Собирались две группы метрик. Первая группа описывает качество игры: количество побед, поражений и партий, завершённых по лимиту ходов. Вторая группа описывает вычислительную стоимость: среднее время выбора хода, глубину поиска и число просмотренных состояний.

Описанный эксперимент проводился на ноутбуке MacBook Pro с процессором Apple M3 Pro, включающим 11 ядер: 5 производительных и 6 энергоэффективных. Объём оперативной памяти составлял 18 ГБ. Операционная система "--- macOS 15.0.

\section{Результаты эксперимента}

Сводная статистика по большому полю приведена в таблице~\ref{tab:summary}. Каждый алгоритм участвовал в трёх попарных сравнениях, поэтому суммарно для него учитывалось 150 партий.

\begin{table}
\centering
\caption{Доля побед алгоритмов на поле $11\times11$}
\label{tab:summary}
\begin{tabular}{lrrrr}
\toprule
Алгоритм & Партий & Победы & Лимит & Доля побед \\
\midrule
AlphaBeta  & 150 & 121 & 2 & 80,67\,\% \\
MonteCarlo & 150 & 88  & 3 & 58,67\,\% \\
MiniMax    & 150 & 84  & 4 & 56,00\,\% \\
Случайный  & 150 & 0   & 5 & 0,00\,\% \\
\bottomrule
\end{tabular}
\end{table}

Наибольшую долю побед в закрытом сравнении на поле $11 \times 11$ показал AlphaBeta.

Для анализа вычислительной стоимости использовались агрегированные средние значения по ходам алгоритмов. Показатели для того же поля $11 \times 11$ приведены в таблице~\ref{tab:cost}.

\begin{table}
\centering
\caption{Вычислительные показатели на поле $11\times11$}
\label{tab:cost}
\resizebox{\textwidth}{!}{%
\begin{tabular}{lrrrr}
\toprule
Показатель & Случайный & MiniMax & MonteCarlo & AlphaBeta \\
\midrule
Длина партии, ходов & 114,87 & 111,43 & 120,27 & 107,17 \\
Время хода, мс & 1,0 & 1578,5 & 6576,2 & 2743,9 \\
Глубина поиска & 1,00 & 3,96 & 108,60 & 6,71 \\
Состояний просмотрено & 4,4 & 2158,1 & 444,7 & 5405,1 \\
Ветвей отброшено & 0,0 & 0,0 & 0,0 & 1140,3 \\
\bottomrule
\end{tabular}
}
\end{table}

MonteCarlo тратит больше всего времени на ход, поскольку выполняет множество симуляций. MiniMax работает быстрее, но не использует отсечения. AlphaBeta требует больше времени, чем MiniMax, зато просматривает более глубокие продолжения и отбрасывает часть ветвей дерева.

\section{Заключение}

В работе рассмотрен модуль программного соперника для игрового приложения <<Коридор>>. Описаны игровая модель, функция оценки позиции и три алгоритма выбора хода. Для сравнения алгоритмов реализован режим автоматических партий, позволяющий собирать как игровые, так и вычислительные метрики.

Эксперимент на поле $11 \times 11$ показал, что среди реализованных стратегий наибольшую долю побед получил AlphaBeta. Он использует отсечения и лучше подходит для более глубокого поиска, тогда как MiniMax проще и быстрее. MonteCarlo оказался конкурентоспособным по качеству игры, но требует существенно большего времени на ход. Код приложения находится в соответствующем репозитории~\cite{github}.

\printbibliography

\end{document}

%%% Local Variables:
%%% mode: latex
%%% TeX-engine: luatex
%%% LaTeX-biblatex-use-Biber: t
%%% TeX-master: t
%%% End:
